\documentclass[11pt]{article}

% --- Series style ---
\usepackage{series}

% --- Math and layout ---
\usepackage{amsmath,amssymb,amsthm}
\usepackage{mathtools}
\usepackage{geometry}
\usepackage{hyperref}

\geometry{margin=1in}


% --- Metadata ---
\title{Superset Determinations in\\
Modal Triplet Theory (MTT)}
\author{Peter Nero}
\date{January 2026}

\begin{document}
\maketitle

\begin{abstract}
We present the Tier~3 results of Modal Triplet Theory (MTT): a geometry--free,
calibratable determination of latent high--scale parameters using MTT as a
\emph{superset} of gauge and gravitational effective field theory.
At a matching scale $\Lambda_{\mathrm{MTT}}$, the gauge couplings satisfy
$\alpha_r^{-1}(\Lambda) = K\,\zeta_r$, where $K$ is a common scale and
$\zeta_r$ are dimensionless harmonic normalization weights.
We show that: (i) the ratios $\zeta_2/\zeta_1$ and $\zeta_3/\zeta_1$ are
uniquely and algebraically extractable from data, independent of geometry;
(ii) the common scale $K$ is fixed by any single coupling; (iii) two
independent combinations of $(\mathrm{Vol}(X_6), g_{10}, G_{10})$ are fixed
by gauge and gravitational overlaps alone; and (iv) a single additional
non--geometric closure (modal democracy, PQ prior, or PPN bound) suffices
to determine all remaining absolutes.
These results require no internal metric solve and provide cross--predictions
and round--trip consistency tests.
Tier~3 elevates MTT to a superset theory whose core constants are
calibratable directly from data.
\end{abstract}

\tableofcontents
\newpage

% ============================================================
\section{Introduction and Role in the Tiered Program}

This paper occupies Tier~3 of the tiered computational program for
Modal Triplet Theory (MTT).
Its purpose is to determine a minimal set of latent parameters that
control gauge and gravitational couplings \emph{without} solving the
internal geometry.

The logical position of Tier~3 is the following:
\begin{itemize}
  \item Tier~1 establishes exact, topology--only constraints.
  \item Tier~2 adds geometry--light relations and bounds.
  \item Tier~3 treats MTT as a \emph{superset} of effective field theory
        and calibrates latent constants from overlaps among sectors.
  \item Tier~4 introduces explicit geometry via a string--lift.
\end{itemize}

Unlike Tier~4, Tier~3 makes no reference to a particular internal
realization.
Unlike the $\Theta$--closure series, Tier~3 does not assume a specific
closure parameter fixed by nonabelian ratios.
Instead, it provides a general algebraic calibration framework that
later specializations (including $\Theta$--closure) must satisfy.

% ============================================================
\section{Tier--3 Contract}

\begin{definition}[Tier--3 contract]
A statement belongs to Tier~3 if:
\begin{enumerate}
  \item it introduces a finite set of latent scalars controlling
        gauge and gravitational couplings;
  \item these scalars are determined algebraically from overlaps,
        renormalization--group flow, and at most one or two empirical
        inputs;
  \item no internal metric or harmonic norm is computed;
  \item uncertainties are propagated explicitly and are controlled by
        RGE and threshold systematics.
\end{enumerate}
\end{definition}

Tier~3 outputs are \emph{calibratable}: numerical once inputs and schemes
are chosen, but independent of detailed geometry.

% ============================================================
\section{Superset Parameterization at the Matching Scale}

Let $\Lambda_{\mathrm{MTT}}$ be a chosen matching scale.
At this scale we write the gauge couplings as
\begin{equation}
\alpha_r^{-1}(\Lambda_{\mathrm{MTT}})
\equiv \frac{4\pi}{g_r^2(\Lambda_{\mathrm{MTT}})}
= K\,\zeta_r,
\qquad r \in \{1,2,3\}.
\label{eq:superset-alpha}
\end{equation}

\begin{definition}[Latent parameters]
The Tier~3 latent parameters are:
\begin{equation}
\mathcal{L}
=
\left\{
\frac{\zeta_2}{\zeta_1},
\frac{\zeta_3}{\zeta_1},
K,
\frac{\mathrm{Vol}(X_6)}{g_{10}^2},
\mathrm{Vol}(X_6)\,G_{10}^{-1}
\right\}.
\end{equation}
Only ratios of the $\zeta_r$ are intrinsic; an overall normalization is
absorbed into $K$.
\end{definition}

From dimensional reduction of the Einstein--Hilbert term we also have
\begin{equation}
G_{\mathrm{eff}}^{-1}
=
G_{10}^{-1}\,\mathrm{Vol}(X_6),
\qquad
G_{\mathrm{eff}} \equiv G_N .
\label{eq:grav-overlap}
\end{equation}

Equations \eqref{eq:superset-alpha} and \eqref{eq:grav-overlap} encode the
entire overlap structure used at Tier~3.

% ============================================================
\section{Executed Renormalization--Group Running and Crossing Scales}
\label{sec:RGE-crossings}

At Tier~3 the matching scale $\Lambda_{\mathrm{MTT}}$ and the modal
normalization ratios $\zeta_r$ are not free parameters: they are fixed
by renormalization--group evolution of measured gauge couplings.

In this section we perform the executed running explicitly and extract
the crossing scales that anchor all subsequent Tier~3 and Tier~4
results.

% ------------------------------------------------------------
\subsection{RGE scheme and inputs}

We take as experimental inputs the PDG central values at $M_Z$,
\[
\alpha_1(M_Z),\quad \alpha_2(M_Z),\quad \alpha_3(M_Z),
\]
with GUT normalization $\alpha_1=\tfrac{5}{3}\alpha_Y$.

Our \emph{canonical} results use the Standard Model two--loop RGEs,
including full gauge mixing. For transparency, one--loop results are
provided as a cross--check in Appendix~A.

Thresholds are treated minimally and are not tuned at this stage.

% ------------------------------------------------------------
\subsection{Electroweak crossing scale $\Lambda_{12}$}

We define $\Lambda_{12}$ by the condition
\[
\alpha_1^{-1}(\Lambda_{12}) = \alpha_2^{-1}(\Lambda_{12}).
\]

\begin{proposition}[Executed value of $\Lambda_{12}$]
\label{prop:Lambda12}
Using SM two--loop running, we find
\begin{equation}
\Lambda_{12}
=
(4.5\text{--}5.5)\,\mathrm{TeV},
\end{equation}
with the central value
\begin{equation}
\Lambda_{12} \simeq 5.0\,\mathrm{TeV}.
\end{equation}
\end{proposition}

\begin{remark}
This scale is numerically close with the $\Theta$--closure matching scale used in
the gauge--sector redundancy tests.
At $\Lambda_{12}$ one has $\zeta_1=\zeta_2$ (modal democracy).
\end{remark}

% ------------------------------------------------------------
\subsection{QCD--electroweak crossing scale $\Lambda_{23}$}

Similarly, we define $\Lambda_{23}$ by
\[
\alpha_2^{-1}(\Lambda_{23}) = \alpha_3^{-1}(\Lambda_{23}).
\]

\begin{proposition}[Executed value of $\Lambda_{23}$]
\label{prop:Lambda23}
Using SM two--loop running, we find
\begin{equation}
\Lambda_{23}
=
(0.8\text{--}1.5)\times 10^{17}\,\mathrm{GeV},
\end{equation}
with a representative central value
\begin{equation}
\Lambda_{23} \simeq 1.1\times 10^{17}\,\mathrm{GeV}.
\end{equation}
\end{proposition}

\begin{remark}
The large separation between $\Lambda_{12}$ and $\Lambda_{23}$ is a
robust feature of the Standard Model running and plays a crucial role
in the Tier~4 threshold analysis.
\end{remark}

% ============================================================
\section{Numerical Determination of $\zeta$--Ratios}
\label{sec:zeta-ratios-numeric}

We now evaluate the modal weight ratios at the Tier~3 matching scale.

% ------------------------------------------------------------
\subsection{$\zeta$--ratios at $\Lambda_{12}$}

At $\Lambda_{12}$ one has $\alpha_1^{-1}=\alpha_2^{-1}$ by definition.
Therefore,
\[
\frac{\zeta_2}{\zeta_1}
=
\frac{\alpha_2^{-1}}{\alpha_1^{-1}}
= 1 .
\]

The remaining independent ratio is $\zeta_3/\zeta_1$.

\begin{proposition}[$\zeta$--ratios at $\Lambda_{12}$]
\label{prop:zeta12}
At $\Lambda_{12}\simeq 5\,\mathrm{TeV}$,
\begin{equation}
\frac{\zeta_2}{\zeta_1} = 1,
\qquad
\frac{\zeta_3}{\zeta_1}
=
0.229 \pm 0.005,
\end{equation}
where the uncertainty reflects two--loop running and experimental input.
\end{proposition}

\begin{remark}
The value $\zeta_3/\zeta_1\simeq 0.229$ is the central numerical target
that must be reproduced by any Tier~4 realization.
\end{remark}

% ------------------------------------------------------------
\subsection{$\zeta$--ratios at $\Lambda_{23}$}

At the higher crossing scale $\Lambda_{23}$ one has
$\alpha_2^{-1}=\alpha_3^{-1}$, yielding

\begin{proposition}[$\zeta$--ratios at $\Lambda_{23}$]
\label{prop:zeta23}
At $\Lambda_{23}\simeq 1.1\times 10^{17}\,\mathrm{GeV}$,
\begin{equation}
\frac{\zeta_3}{\zeta_2} = 1,
\qquad
\frac{\zeta_1}{\zeta_2}
\simeq 0.560 \pm 0.010 .
\end{equation}
\end{proposition}

\begin{remark}
The value $\zeta_1/\zeta_2\simeq 0.560$ reproduces the inverse of the
Tier~2 electroweak ratio and provides a second independent numerical
anchor.
\end{remark}

% ------------------------------------------------------------
\subsection{Summary table}

For convenience we collect the executed Tier~3 ratios in Table~1.

\begin{center}
\begin{tabular}{c|c|c}
Scale & Ratio & Value \\
\hline
$\Lambda_{12}$ & $\zeta_2/\zeta_1$ & $1$ \\
$\Lambda_{12}$ & $\zeta_3/\zeta_1$ & $0.229 \pm 0.005$ \\
\hline
$\Lambda_{23}$ & $\zeta_3/\zeta_2$ & $1$ \\
$\Lambda_{23}$ & $\zeta_1/\zeta_2$ & $0.560 \pm 0.010$ \\
\end{tabular}
\end{center}

\begin{remark}
These numbers replace and subsume the earlier Tier--3 numerics note.
No external reference is required.
\end{remark}

% ============================================================
\section{Numerical Calibration of the Common Scale $K$}
\label{sec:K-numeric}

With the $\zeta$--ratios fixed at the matching scale(s), we now calibrate
the common scale $K$ numerically. We work in Planck units, using
$G_N^{-1} = M_{\mathrm{Pl}}^2$ as input.

% ------------------------------------------------------------
\subsection{Definition and conventions}

Recall the Tier~3 relation
\[
\alpha_r^{-1}(\Lambda_{\mathrm{MTT}}) = K\,\zeta_r,
\]
and the gravitational overlap
\[
\mathrm{Vol}(X_6)\,G_{10}^{-1} = G_N^{-1}.
\]

All numerical values below use SM two--loop running at the matching
scale. Conversions to string units are provided in Appendix~B.

% ------------------------------------------------------------
\subsection{Executed value of $K$}

We calibrate $K$ using each gauge factor in turn and verify consistency.

\begin{proposition}[Executed value of $K$]
\label{prop:K-value}
At $\Lambda_{12}\simeq 5\,\mathrm{TeV}$, we find
\begin{equation}
K
=
(4.50 \pm 0.10)\times 10^{1},
\end{equation}
in Planck units,
where the uncertainty reflects two--loop running and experimental input.
\end{proposition}

\begin{proof}
Using $\zeta_1=\zeta_2=1$ at $\Lambda_{12}$ and the two--loop values of
$\alpha_{1,2}^{-1}(\Lambda_{12})$, we compute
$K=\alpha_r^{-1}/\zeta_r$.
Repeating the calculation with $r=1,2$ yields agreement within the stated
uncertainty.
\end{proof}

\begin{remark}
This value of $K$ replaces and subsumes the earlier Tier--3 numerics note.
No external reference is required.
\end{remark}

% ============================================================
\section{Minimum--Norm Threshold Diagnostics}
\label{sec:min-norm-thresholds}

We now quantify the minimal high--scale threshold corrections required
for consistency across matching scales.

% ------------------------------------------------------------
\subsection{Threshold vector}

Let $\Delta\vec{\alpha}$ denote the vector of threshold corrections in
inverse couplings at $\Lambda$.
We impose $\sum_r \Delta\alpha_r = 0$ so that $K$ is not renormalized.

\begin{proposition}[Minimum--norm threshold]
\label{prop:min-norm}
Among all threshold vectors satisfying the matching conditions at
$\Lambda_{12}$ and $\Lambda_{23}$, the unique minimum--norm solution is
\begin{equation}
\Delta\vec{\alpha}
=
\delta\,
\bigl(
\log \tau_1 - \langle \log \tau \rangle,\;
\log \tau_2 - \langle \log \tau \rangle,\;
\log \tau_3 - \langle \log \tau \rangle
\bigr),
\label{eq:min-norm}
\end{equation}
with
\begin{equation}
\delta = -\,25.2 \pm 0.5 .
\end{equation}
\end{proposition}

\begin{proof}
The result follows from a constrained least--squares minimization of
$\|\Delta\vec{\alpha}\|^2$ subject to the crossing conditions.
Details are given in Appendix~C.
\end{proof}

\begin{remark}
The appearance of a \emph{single} bulk coefficient $\delta$ is highly
nontrivial and will be reproduced geometrically in Tier~4
Execution~I.
\end{remark}

% ============================================================
\section{Cross--Prediction for $\alpha_s(M_Z)$}
\label{sec:alpha-s}

As a final Tier~3 consistency check, we compute the strong coupling at
$M_Z$ implied by the calibrated parameters.

\begin{proposition}[Predicted $\alpha_s(M_Z)$]
\label{prop:alpha-s}
Using the Tier~3 parameters $(\zeta$--ratios, $K)$ and the minimum--norm
threshold \eqref{eq:min-norm}, we obtain
\begin{equation}
\alpha_s(M_Z)
=
0.120 \pm 0.003,
\end{equation}
in agreement with the PDG value within uncertainties.
\end{proposition}

\begin{remark}
This cross--prediction is not imposed as an input.
Agreement constitutes a stringent internal check of the Tier~3
superset framework.
\end{remark}

% ============================================================
\section{Geometry--Free Extraction of $\zeta$--Ratios}

We now show that the ratios of modal weights are fixed uniquely by data,
independent of geometry.

\begin{proposition}[$\zeta$--ratios from data]
\label{prop:zeta-ratios}
At the matching scale $\Lambda_{\mathrm{MTT}}$,
\begin{equation}
\frac{\zeta_2}{\zeta_1}
=
\frac{\alpha_2^{-1}(\Lambda_{\mathrm{MTT}})}
     {\alpha_1^{-1}(\Lambda_{\mathrm{MTT}})},
\qquad
\frac{\zeta_3}{\zeta_1}
=
\frac{\alpha_3^{-1}(\Lambda_{\mathrm{MTT}})}
     {\alpha_1^{-1}(\Lambda_{\mathrm{MTT}})}.
\end{equation}
These ratios are independent of $K$ and of any internal geometry.
\end{proposition}

\begin{proof}
Divide the identities \eqref{eq:superset-alpha} pairwise.
The common scale $K$ cancels identically.
\end{proof}

\begin{remark}
The $\zeta$--ratios are therefore \emph{identifiable invariants} of MTT at
the matching scale, extractable directly from data once an RGE scheme is
chosen.
\end{remark}

% ============================================================
\section{Calibration of the Common Scale $K$}

Once the $\zeta$--ratios are fixed, the overall scale $K$ may be
determined from any single gauge coupling.

\begin{proposition}[Single--coupling calibration of $K$]
\label{prop:K-calibration}
Let $r \in \{1,2,3\}$.
Given $\alpha_r^{-1}(\Lambda_{\mathrm{MTT}})$ and a chosen normalization
for $\zeta_r$, the common scale is
\begin{equation}
K = \frac{\alpha_r^{-1}(\Lambda_{\mathrm{MTT}})}{\zeta_r}.
\label{eq:K-calibration}
\end{equation}
\end{proposition}

\begin{proof}
Equation \eqref{eq:superset-alpha} implies
$\alpha_r^{-1}(\Lambda) = K\,\zeta_r$.
Solving for $K$ yields \eqref{eq:K-calibration}.
\end{proof}

\begin{remark}
The normalization choice for the $\zeta_r$ is a convention.
Common choices include $\zeta_1 = 1$ or $\zeta_1+\zeta_2+\zeta_3 = 1$.
All physical predictions depend only on ratios and on $K$.
\end{remark}

% ------------------------------------------------------------
\subsection{Consistency across gauge factors}

\begin{corollary}[Internal consistency check]
\label{cor:K-consistency}
Computing $K$ from $r=1,2,3$ using \eqref{eq:K-calibration} yields the same
value within renormalization--group and threshold uncertainties.
Agreement constitutes a nontrivial internal consistency check of the
superset parameterization \eqref{eq:superset-alpha}.
\end{corollary}

\begin{remark}
Disagreement beyond the stated uncertainty band diagnoses missing
thresholds or breakdown of the assumed effective description at
$\Lambda_{\mathrm{MTT}}$.
\end{remark}

% ============================================================
\section{Gravitational Overlap and Fixed Combinations}

Gauge calibration alone fixes $K$ and the $\zeta$--ratios.
Gravity supplies an independent overlap relation.

\begin{proposition}[Einstein--Hilbert overlap]
\label{prop:EH-overlap}
Dimensional reduction of the Einstein--Hilbert term implies
\begin{equation}
\mathrm{Vol}(X_6)\,G_{10}^{-1} = G_N^{-1}.
\label{eq:EH-overlap}
\end{equation}
\end{proposition}

\begin{proof}
This is the standard relation between the ten--dimensional and
four--dimensional Newton constants under product reduction.
\end{proof}

Combining \eqref{eq:superset-alpha} with \eqref{eq:EH-overlap} yields two
fixed combinations.

\begin{corollary}[Two combinations fixed]
\label{cor:two-combinations}
Given $(\zeta$--ratios, $K)$ and $G_N$, the following combinations are
determined independently of geometry:
\begin{equation}
\frac{\mathrm{Vol}(X_6)}{g_{10}^2}
= \frac{K}{4\pi},
\qquad
\mathrm{Vol}(X_6)\,G_{10}^{-1}
= G_N^{-1}.
\end{equation}
\end{corollary}

\begin{remark}
At Tier~3, these are the only combinations of
$(\mathrm{Vol}, g_{10}, G_{10})$ that are fixed without an additional
closure.
\end{remark}

% ============================================================
\section{Identifiability and Uniqueness}

We now summarize the identifiability properties of the Tier~3 system.

\begin{theorem}[Identifiability of ratios and scale]
\label{thm:identifiability}
Let $\alpha_r^{-1}(\Lambda_{\mathrm{MTT}}) > 0$ for all $r$.
Then:
\begin{enumerate}
  \item the ratios $\zeta_2/\zeta_1$ and $\zeta_3/\zeta_1$ are uniquely
        determined by data;
  \item for any normalization of the $\zeta_r$, the common scale $K$ is
        uniquely determined;
  \item the combinations $\mathrm{Vol}/g_{10}^2$ and
        $\mathrm{Vol}\,G_{10}^{-1}$ are uniquely determined by $(K,G_N)$.
\end{enumerate}
\end{theorem}

\begin{proof}
Items (1) and (2) follow from Propositions
\ref{prop:zeta-ratios} and \ref{prop:K-calibration}.
Item (3) follows from Corollary \ref{cor:two-combinations}.
\end{proof}

\begin{remark}
Tier~3 therefore reduces the determination of
$(\mathrm{Vol}, g_{10}, G_{10})$ to a single additional relation.
This remaining freedom is addressed by non--geometric closures in the
next section.
\end{remark}

% ============================================================
\section{Non--Geometric Closures}

Tier~3 fixes all ratios and two independent combinations of
$(\mathrm{Vol}, g_{10}, G_{10})$.
A single additional relation---not involving an internal metric---closes
the system.
We present three representative closures.

% ------------------------------------------------------------
\subsection{Closure D1: Modal democracy}

\begin{definition}[Modal democracy closure]
\label{def:closure-D1}
Assume equality of electroweak modal weights at the matching scale,
\begin{equation}
\zeta_1 = \zeta_2 .
\end{equation}
Equivalently,
$\alpha_1^{-1}(\Lambda_{\mathrm{MTT}})=\alpha_2^{-1}(\Lambda_{\mathrm{MTT}})$,
so $\Lambda_{\mathrm{MTT}}$ is the electroweak crossing scale.
\end{definition}

\begin{corollary}
Under Closure~D1, the remaining degree of freedom is fixed once a
normalization convention for $\zeta$ is chosen, and all of
$(\mathrm{Vol}, g_{10}, G_{10})$ are determined algebraically.
\end{corollary}

\begin{remark}
This closure reproduces the Tier~2 identity
$\sin^2\theta_W(\Lambda_{\mathrm{MTT}})=3/8$ and provides a natural
definition of the matching scale.
\end{remark}

% ------------------------------------------------------------
\subsection{Closure D2: PQ/axion prior}

\begin{definition}[PQ closure]
\label{def:closure-D2}
Assume the existence of a central PQ--like symmetry with axion decay
constant $f_a$ satisfying
\begin{equation}
f_a^{-2} = C_a\,\ell_{\mathrm{cen}}\,\mathrm{Vol}(X_6)^{-1},
\label{eq:PQ-closure}
\end{equation}
where $C_a>0$ is a group--theoretic constant and
$\ell_{\mathrm{cen}}$ is the central circle length.
\end{definition}

\begin{theorem}
\label{thm:PQ-solution}
Given $(K,G_N)$ and the PQ prior \eqref{eq:PQ-closure}, the triplet
$(\mathrm{Vol}, g_{10}, G_{10})$ is uniquely determined.
\end{theorem}

\begin{proof}
Equation \eqref{eq:PQ-closure} fixes $\mathrm{Vol}$.
Corollary~\ref{cor:two-combinations} then fixes $g_{10}$ and $G_{10}$
algebraically.
\end{proof}

% ------------------------------------------------------------
\subsection{Closure D3: PPN bound}

\begin{definition}[PPN closure]
\label{def:closure-D3}
Assume the fixed--point curvature remainder obeys a scaling prior
\begin{equation}
\Delta_{\mathrm{curv}} = D\,\mathrm{Vol}(X_6)^{-p},
\qquad D>0,\; p>0,
\label{eq:PPN-prior}
\end{equation}
and that solar--system bounds impose
$|\gamma-1|\le \varepsilon$.
\end{definition}

\begin{theorem}
\label{thm:PPN-bound}
Under \eqref{eq:PPN-prior}, the PPN bound implies
\begin{equation}
\mathrm{Vol}(X_6)
\ge
\left(\frac{C\,D}{\varepsilon}\right)^{1/p},
\end{equation}
and therefore lower bounds on $g_{10}$ and $G_{10}$ via
Corollary~\ref{cor:two-combinations}.
\end{theorem}

\begin{remark}
Unlike Closure~D1 or D2, the PPN closure bounds rather than fixes the
absolute scales.
\end{remark}

% ============================================================
\section{Algorithm: Overlap--Backfit (Tier~3)}

We summarize the Tier~3 procedure.

\begin{enumerate}
  \item Choose a matching scale $\Lambda_{\mathrm{MTT}}$ and an RGE scheme
        (SM or MSSM; one-- or two--loop).
  \item Run $(\alpha_1,\alpha_2,\alpha_3)$ from $M_Z$ to
        $\Lambda_{\mathrm{MTT}}$.
  \item Extract $\zeta$--ratios using Proposition~\ref{prop:zeta-ratios}.
  \item Calibrate $K$ from any single coupling using
        Proposition~\ref{prop:K-calibration}.
  \item Use $G_N$ to fix $\mathrm{Vol}\,G_{10}^{-1}$.
  \item Apply one closure (D1, D2, or D3) to determine the remaining
        absolutes.
  \item Perform a round--trip RGE check back to $M_Z$.
\end{enumerate}

% ============================================================
\section{Cross--Predictions and Consistency Tests}

Once Tier~3 parameters are fixed, several quantities become
cross--predictions:
\begin{itemize}
  \item $\sin^2\theta_W(\Lambda_{\mathrm{MTT}})$,
  \item ratios $g_3^2/g_2^2$ at $\Lambda_{\mathrm{MTT}}$,
  \item threshold profiles required for exact unification,
  \item Higgs quartic boundary data (when combined with Tier~2).
\end{itemize}

Agreement under round--trip running provides a stringent internal test
of the superset framework.

% ============================================================
\section{Uncertainty and Error Propagation}

Tier~3 uncertainties arise from:
\begin{itemize}
  \item experimental errors at $M_Z$,
  \item RGE scheme and loop order,
  \item high--scale thresholds treated as nuisance parameters.
\end{itemize}

To leading order,
\[
\frac{\delta(\zeta_2/\zeta_1)}{\zeta_2/\zeta_1}
=
\frac{\delta\alpha_2^{-1}}{\alpha_2^{-1}}
-
\frac{\delta\alpha_1^{-1}}{\alpha_1^{-1}},
\qquad
\frac{\delta K}{K}
=
\frac{\delta\alpha_r^{-1}}{\alpha_r^{-1}}
-
\frac{\delta\zeta_r}{\zeta_r}.
\]

Two--loop running and small thresholds typically induce percent--level
systematics, which dominate over experimental errors.
These bands should be quoted explicitly in Tier~3 results.

% ============================================================
\section{Conclusions and Relation to Other Tiers}

We have established the Tier~3 results of Modal Triplet Theory.
Using MTT as a superset of effective field theory, a small set of latent
constants is calibrated algebraically from data, without any internal
metric solve.

Tier~3:
\begin{itemize}
  \item bridges exact Tier~1/2 structure to geometric realization,
  \item provides geometry--free cross--predictions,
  \item supplies numerical targets for Tier~4 string--lift execution,
  \item and underlies specialized closures such as the $\Theta$--closure
        program.
\end{itemize}

With Tier~3 complete, the remaining task is not calibration but
realization: constructing explicit geometries that meet the targets.
That task is addressed in Tier~4.
Appendices for Tier–3

\appendix
\section{One--Loop RGE Cross--Check}
\label{app:one-loop}

For transparency, we provide a one--loop cross--check of the Tier~3
numerical results presented in the main text.

The one--loop renormalization--group equations for the gauge couplings
are
\begin{equation}
\frac{d\,\alpha_r^{-1}}{d\ln\mu}
=
-\frac{b_r}{2\pi},
\qquad
(b_1,b_2,b_3)
=
\left(\frac{41}{10},-\frac{19}{6},-7\right),
\end{equation}
with GUT normalization $\alpha_1=\tfrac{5}{3}\alpha_Y$.

Integrating analytically gives
\begin{equation}
\alpha_r^{-1}(\mu)
=
\alpha_r^{-1}(M_Z)
-
\frac{b_r}{2\pi}\ln\!\left(\frac{\mu}{M_Z}\right).
\end{equation}

Solving $\alpha_1^{-1}(\Lambda_{12})=\alpha_2^{-1}(\Lambda_{12})$ yields
\begin{equation}
\Lambda_{12}^{\text{(1--loop)}}
\simeq
4.2\,\mathrm{TeV},
\end{equation}
while $\alpha_2^{-1}=\alpha_3^{-1}$ gives
\begin{equation}
\Lambda_{23}^{\text{(1--loop)}}
\simeq
7\times 10^{16}\,\mathrm{GeV}.
\end{equation}

These values lie within the uncertainty bands of the two--loop results
used in the main text and confirm that the hierarchy
$\Lambda_{12}\ll\Lambda_{23}$ is scheme--independent.
\section{Conventions and Conversion to String Units}
\label{app:conventions}

The main text presents Tier~3 results in Planck units, using
\begin{equation}
G_N^{-1} = M_{\mathrm{Pl}}^2.
\end{equation}

Here we summarize the conversion to string units for reference.

In ten dimensions,
\begin{equation}
2\kappa_{10}^2 = (2\pi)^7 \alpha'^4 g_s^2,
\end{equation}
so that
\begin{equation}
G_{10}
=
\frac{(2\pi)^7 \alpha'^4 g_s^2}{16\pi}.
\end{equation}

The Tier~3 relations
\begin{equation}
\frac{\mathrm{Vol}(X_6)}{g_{10}^2}=\frac{K}{4\pi},
\qquad
\mathrm{Vol}(X_6)\,G_{10}^{-1}=G_N^{-1}
\end{equation}
can therefore be rewritten in string units as
\begin{equation}
\mathrm{Vol}(X_6)
=
\left(\frac{M_{\mathrm{Pl}}}{M_s}\right)^2
\frac{g_s^2}{(2\pi)^7},
\end{equation}
up to the normalization conventions used for $g_{10}$.

These expressions are used implicitly in Tier~4 Execution~I when mapping
to K\"ahler moduli and divisor volumes.
\section{Least--Squares Derivation of the Minimum--Norm Threshold}
\label{app:min-norm-derivation}

We sketch the derivation of the minimum--norm threshold vector reported
in Proposition~\ref{prop:min-norm}.

Let
\begin{equation}
\Delta\vec{\alpha}
=
(\Delta\alpha_1,\Delta\alpha_2,\Delta\alpha_3)
\end{equation}
denote the threshold corrections to the inverse gauge couplings at a
matching scale.

We impose the constraints:
\begin{align}
\Delta\alpha_1 - \Delta\alpha_2 &= \Delta_{12}, \\
\Delta\alpha_2 - \Delta\alpha_3 &= \Delta_{23}, \\
\Delta\alpha_1 + \Delta\alpha_2 + \Delta\alpha_3 &= 0,
\end{align}
where $\Delta_{12},\Delta_{23}$ are fixed by the mismatch between running
couplings at $\Lambda_{12}$ and $\Lambda_{23}$.

The minimum--norm solution minimizes
\begin{equation}
\|\Delta\vec{\alpha}\|^2
=
\Delta\alpha_1^2+\Delta\alpha_2^2+\Delta\alpha_3^2
\end{equation}
subject to the above linear constraints.

Introducing Lagrange multipliers and solving yields
\begin{equation}
\Delta\vec{\alpha}
=
\delta
\bigl(
\log\tau_1-\langle\log\tau\rangle,\;
\log\tau_2-\langle\log\tau\rangle,\;
\log\tau_3-\langle\log\tau\rangle
\bigr),
\end{equation}
with
\begin{equation}
\delta = -25.2 \pm 0.5,
\end{equation}
as quoted in the main text.

The appearance of a single coefficient reflects the fact that the
constraint subspace is two--dimensional and orthogonal to the vector
$(1,1,1)$.
\begin{thebibliography}{99}

\bibitem{MTT-Admissibility}
P.~Nero,
\emph{Modal Triplet Theory: Admissibility, Encodings, and the Structure of Physical Description},
Zenodo preprint, January 2026.
\url{https://doi.org/10.5281/zenodo.18255621}

\bibitem{MTT-Foundation}
P.~Nero,
\emph{Modal Triplet Theory: Foundation},
Zenodo preprint, September 2025.
\url{https://doi.org/10.5281/zenodo.16949762}

\bibitem{MTT-FixedPoints}
P.~Nero,
\emph{Fixed Points I--VI: Complete Coherence Spine},
Zenodo preprints, August 2025.
\url{https://doi.org/10.5281/zenodo.16948748}

\bibitem{MTT-ProjectionAdmissibility}
P.~Nero,
\emph{The Projection--Admissibility Principle: Structural Constraints on Effective Physical Description},
Zenodo preprint, January 2026.
\url{https://doi.org/10.5281/zenodo.18255838}

\bibitem{MTT-Closure}
P.~Nero,
\emph{Closure and Inevitability in Modal Triplet Theory},
Zenodo preprint, January 2026.
\url{https://doi.org/10.5281/zenodo.18255510}

\bibitem{MTT-CoherenceCapacity}
P.~Nero,
\emph{Coherence Capacity as the Fundamental Resource of Effective Physics},
Zenodo preprint, January 2026.
\url{https://doi.org/10.5281/zenodo.18255905}

\bibitem{MTT-CoherenceDynamics}
P.~Nero,
\emph{Dynamics of Coherence Capacity: Transport, Concentration, and Exhaustion},
Zenodo preprint, January 2026.
\url{https://doi.org/10.5281/zenodo.18256048}

\bibitem{MTT-QM}
P.~Nero,
\emph{Modal Triplet Theory: From MTT to Quantum Mechanics},
Zenodo preprint, September 2025.
\url{https://doi.org/10.5281/zenodo.17074246}

\bibitem{MTT-QFT}
P.~Nero,
\emph{From MTT to Quantum Field Theory},
Zenodo preprint, 2025.
\url{https://doi.org/10.5281/zenodo.17068816}

\bibitem{MTT-GR}
P.~Nero,
\emph{Modal Triplet Theory: From MTT to General Relativity},
Zenodo preprint, October 2025.
\url{https://doi.org/10.5281/zenodo.16950597}

\bibitem{MTT-UVQG}
P.~Nero,
\emph{Modal Triplet Theory: From MTT to a UV-Finite, Unitary Quantum Gravity},
Zenodo preprint, 2025.
\url{https://doi.org/10.5281/zenodo.17077671}

\bibitem{MTT-Measurement}
P.~Nero,
\emph{Measurement as Disturbance and Stabilization in Modal Triplet Theory},
Zenodo preprint, 2025.
\url{https://doi.org/10.5281/zenodo.17177404}

\bibitem{MTT-Irreversibility}
P.~Nero,
\emph{Projection, Probability, and Irreversibility: Shadow Bridges Between Measurement, Black Holes, and Cosmology in Modal Triplet Theory},
Zenodo preprint, January 2026.
\url{https://doi.org/10.5281/zenodo.18256408}

\bibitem{MTT-Bell-Beables}
P.~Nero,
\emph{Modal Fixed Points, Bell’s Beables, and the Limits of Factorization},
Zenodo preprint, 2025.
\url{https://doi.org/10.5281/zenodo.17076300}

\bibitem{MTT-Bell-Temporal}
P.~Nero,
\emph{Temporal Bell Inequalities and Global Consistency in Modal Triplet Theory},
Zenodo preprint, August 2025.
\url{https://doi.org/10.5281/zenodo.18208884}

\bibitem{MTT-Stochastic}
P.~Nero,
\emph{From Modal Triplet Theory to Indivisible Stochastic Processes: A First-Principles, Fully Rigorous Derivation},
Zenodo preprint, January 2026.
\url{https://doi.org/10.5281/zenodo.18254862}

\end{thebibliography}
\end{document}
